% Preamble requirements:
% \usepackage{graphicx}
% \usepackage{tikz}
% \usetikzlibrary{arrows.meta,positioning}
% This file includes a native TikZ figure: no SVG conversion is required.

\subsection{Implementation and Reproducibility}
\label{sec:implementation}

\paragraph{Project organization.}
The application is implemented in Python and separates reusable source code, datasets, interactive experiments, trained models, evaluation outputs, and tests. The installable package resides in \texttt{src/toxic\_comments/}. The \texttt{data/} directory contains raw datasets, processed datasets, reusable fold files, and external pretrained embeddings. The \texttt{notebooks/} directory provides workflows for exploratory analysis, cleaning, embedding comparisons, and neural-model experiments. Fitted model artifacts are stored in the root-level \texttt{models/} directory, whereas model implementations reside in \texttt{src/toxic\_comments/models/}. Metrics, summaries, logs, and figures are stored in \texttt{results/}; automated tests and report materials are maintained in \texttt{tests/} and \texttt{report/}, respectively. Figure~\ref{fig:software-architecture} summarizes the execution paths and dataset-access abstraction.

\begin{figure*}[t]
\centering
\resizebox{\textwidth}{!}{%
\begin{tikzpicture}[
  box/.style={draw, rounded corners=2pt, align=center, font=\small, minimum height=12mm, text width=47mm},
  wide/.style={box,text width=94mm},
  arr/.style={-{Stealth[length=2mm]},semithick},
  impl/.style={arr,dashed},
  lab/.style={font=\scriptsize,fill=white,inner sep=2pt}
]
\node[wide] (entry) at (6,0) {\textbf{Entry points}\\Jupyter notebooks and command-line interface};
\node[box] (experiment) at (3,-2) {\texttt{experiment.py}\\\texttt{run\_experiment(repository, ...)}};
\node[box] (train) at (10,-2) {\texttt{train.py}\\holdout / folds / full};
\node[box,draw=blue!50!black,fill=blue!4] (repo) at (0,-4.5) {\textbf{DatasetRepository}\ (Protocol)\\\texttt{load() -> DataFrame}\\\texttt{save(data)}};
\node[wide] (shared) at (8.5,-4.5) {\textbf{Shared package components}\\cleaning, splits, embeddings, model registry,\\model implementations, evaluation};
\node[box] (memory) at (-1.5,-7) {\texttt{InMemoryDatasetRepository}\\DataFrame copies};
\node[box] (csv) at (4,-7) {\texttt{CsvFileDatasetRepository}\\CSV read / write};
\node[box] (out) at (11,-7) {\textbf{Filesystem outputs}\\\texttt{models/}: fitted artifacts\\\texttt{results/}: metrics};
\node[box] (data) at (4,-9.5) {\textbf{Dataset files in data/}\\raw / processed / k\_fold};
\draw[arr] (entry.south) -- ++(0,-.35) -| node[lab,pos=.75]{invoke} (experiment.north);
\draw[arr] (entry.south) -- ++(0,-.35) -| node[lab,pos=.75]{invoke} (train.north);
\draw[arr] (experiment.west) -| node[lab,pos=.65]{calls load()} (repo.north);
\draw[arr] (experiment.south) -- ++(0,-.65) -| node[lab,pos=.75]{uses} (shared.north west);
\draw[arr] (train.south) -- node[lab]{uses} (shared.north east);
\draw[impl] (memory.north) -- ++(0,.5) -| node[lab,pos=.6]{implements} (repo.south west);
\draw[impl] (csv.north) -- ++(0,.6) -| node[lab,pos=.65]{implements} (repo.south east);
\draw[arr] (shared.south) -- ++(0,-.5) -| node[lab,pos=.7]{workflows persist} (out.north);
\draw[semithick,{Stealth[length=2mm]}-{Stealth[length=2mm]}] (csv.south) -- node[lab]{read / write} (data.north);
\draw[arr] (train.east) -- ++(1.4,0) |- node[lab,pos=.75]{direct CSV access in train.py} (data.east);
\node[align=center,font=\small,text width=150mm] at (5.5,-11) {External resources: \texttt{data/embeddings/} supplies pretrained vectors.\\Quality and provenance: \texttt{tests/}, dependency declarations, Git, and notebook fingerprints.};
\end{tikzpicture}%
}
\caption{Software architecture and interchangeable dataset backends. The repository-based experiment runner consumes the same DataFrame contract for memory and CSV storage. Dashed arrows indicate interface implementation; solid arrows indicate calls or access. The training workflow also supports direct CSV access. Models and results remain filesystem-based artifacts.}
\label{fig:software-architecture}
\end{figure*}

\paragraph{Application components and execution.}
Within the package, \texttt{config.py} centralizes label names, text columns, and directory paths; \texttt{cleaning.py} implements text transformations; and \texttt{splits.py} generates holdout or k-fold partitions and exports reusable fold CSV files. The \texttt{embeddings/} subpackage provides pretrained-vector loaders and pooling operators, while \texttt{build\_models} in the model registry constructs named estimators. Shared evaluation utilities extract continuous scores through \texttt{predict\_proba} or \texttt{decision\_function} and compute aggregate and per-label metrics.

The repository-based \texttt{run\_experiment} workflow prepares data, selects models, generates shared fold indices, and invokes \texttt{cross\_validate\_model}, which clones and fits an estimator for each fold. A complementary workflow in \texttt{train.py} provides \texttt{train\_model}, \texttt{train\_holdout}, and \texttt{train\_from\_folds}, together with a full-dataset training mode. Both workflows reuse model construction and evaluation utilities, but differ in orchestration and persistence. The \texttt{predict.py} and \texttt{submission.py} modules support inference from saved estimators and generation of six-label submission files, respectively.

\paragraph{Interchangeable dataset backends.}
The \texttt{DatasetRepository} protocol defines \texttt{load()}, returning a pandas DataFrame, and \texttt{save(data)}, storing a DataFrame. \texttt{InMemoryDatasetRepository} maintains data in memory and copies it on loading and saving, while \texttt{CsvFileDatasetRepository} implements the same operations using filesystem CSV files. The repository is supplied to \texttt{run\_experiment} as an argument. Switching backends therefore changes only the repository instance at the call site; the DataFrame representation and subsequent validation, preprocessing, training, and evaluation logic remain unchanged.

This abstraction addresses the requirement for interchangeable memory and file-based dataset management. Its scope is dataset access in the repository-based workflow: other entry points can read CSV files directly, and model and metric persistence remain filesystem-based. The runner currently consumes the \texttt{load()} operation; both backends additionally expose \texttt{save(data)} for dataset updates.

\paragraph{Persistence and experiment traceability.}
The \texttt{save\_kfold\_datasets} utility writes reusable \texttt{fold\_N/train.csv} and \texttt{fold\_N/test.csv} partitions. Training utilities serialize fitted estimators with \texttt{joblib}, including fold-specific models, and save aggregate and per-label metrics separately. The repository-based cross-validation workflow writes completed fold metrics incrementally, preserving those records if later computation is interrupted. This does not automatically resume unfinished training. Separate output directories are required to retain distinct runs because some workflows replace existing outputs.

Settings are supplied through function arguments, command-line options, and notebook cells. The embedding-comparison notebook additionally computes experiment fingerprints from package source code, supplied settings, and the relevant training and validation data. Result records carry these fingerprints, allowing completed configurations to be recognized during resumed experiment sweeps. Source-change checks require a kernel restart when imported code may be stale. This mechanism is specific to that notebook and does not capture the complete software environment or all external resource versions.

\paragraph{Testing and reproducibility.}
The pytest suite covers cleaning, split generation, score extraction, metric aggregation, model registration, embedding and neural components, threshold tuning, training, prediction, and submission formatting. Existing tests check reproducible split generation, holdout disjointness, fold-file creation, and model persistence using small datasets and temporary directories. These tests support regression detection; they do not substitute for predictive evaluation or establish that the complete suite passes in every environment.

Random-state settings support repeatable partitioning and stochastic model components. Dependencies are declared in \texttt{pyproject.toml} and \texttt{requirements.txt}, including optional groups for additional model families. These declarations use version constraints rather than a fully locked environment. Consequently, exact reproduction additionally requires the resolved dependency versions, hardware details, pretrained resource versions, data partitions, and model settings. Fixed seeds alone do not guarantee identical neural-network outputs across environments. Git tracks source changes and development contributions; reported runs should be associated with a commit and their execution configuration so that saved metrics and model artifacts can be traced to the corresponding implementation.
